\section{Proposed practitioner evaluation}\label{app:practitioner-protocol}

The artificial evaluation provides technical evidence on integrity, retained evidence, execution controls, minimized disclosure and cost (DR1--DR5) but not on utility for the people who must act on the evidence. Following FEDS \citep{venable2016}, we plan a naturalistic episode. \todo{Lakshman and Prof.\ Neouchi: approve, adapt, or replace the following design, and obtain UT Dallas IRB determination before recruiting.}
\begin{itemize}
\item \emph{Participants:} 12--20 participants in two roles, software developers who build LLM features and auditors or compliance analysts.
\item \emph{Tasks:} (a) developers instrument an existing LLM feature with receipts, measuring time to adoption and code changes (DR5); (b) auditors receive a set of runs, some tampered and some with policy violations, as either conventional logs or receipts with the command-line tool, and must identify which runs are trustworthy and which violated policy (DR1, DR3).
\item \emph{Design:} within-subject comparison of logs and receipts with counterbalanced order.
\item \emph{Measures:} task accuracy and time, perceived usefulness and ease of use, and semi-structured interview data on fit with audit procedures and frameworks such as ISO/IEC 42001 \citep{iso42001}.
\item \emph{Analysis:} paired tests for accuracy and time; thematic analysis of interviews to refine the design principles.
\end{itemize}

